% =====================================================================================
%  Chapter 2 — Background and Related Work
%  Thesis: Comparison of Control Strategies for Interaction-Force Aware Multirotor Teams
%
%  DRAFT, 23 August 2026 (writing-track item W.5).
%
%  CITATION DISCIPLINE — read before editing.
%  THIS FILE IS THE SOURCE OF TRUTH for every factual claim about prior work. Verify a claim
%  against the paper itself. Do NOT verify it against, or import it from, any summary, matrix,
%  ledger or review brief in docs/ — those are stale working notes, archived on 2026-08-27 to
%  docs/_archive/citation_notes/, and where they disagree with this file, THIS FILE WINS.
%  Three misattributions in this project came from treating such a note as a source.
%
%  A number belongs here only with what it is a number OF: which axis, which training set,
%  which platform, which experiment, and whether it comes from the abstract or the body.
%  That distinction — not arithmetic — is what every correction pass has turned on.
%
%  Bibliography: docs/references.bib (generated from the arXiv API).
%
%  The Lee et al. citation is RESOLVED: CDC 2010 is the primary reference; its equations were
%  verified against the companion arXiv paper. The remaining \todo{} in §2.9 is a PROVENANCE
%  NOTE recording which audits support the gap claim, not outstanding work.
%
%  Citation baseline: 2026-08-27 (see docs/thesis/CITATION_BASELINE.md).
%  After that date: only NEW cites or EDITED prior-work sentences need re-checking.
% =====================================================================================

\chapter{Background and Related Work}
\label{ch:related}

Multirotors flying in close proximity disturb one another aerodynamically. The disturbance is
large enough to destabilise a vehicle, and conventional practice avoids it by keeping the vehicles
far apart: a safety distance of roughly \SI{60}{\centi\metre} has been used for a quadrotor only
\SI{9}{\centi\metre} across~\cite{shi2020neuralswarm}. Removing that margin is what makes dense
formation flight possible, and doing so requires a controller that accounts for the interaction
rather than avoiding it.

This chapter surveys how that has been attempted. It is organised as an argument rather than a
catalogue. Section~\ref{sec:rw-physics} establishes what the disturbance is; \ref{sec:rw-reactive}
and \ref{sec:rw-predictive} describe the two sources of information a controller can have about it,
namely online measurement and offline prediction; \ref{sec:rw-inject} shows that a prediction is
only useful once it enters a control law, and that how it enters differs by controller family;
\ref{sec:rw-rl} covers the alternative of learning the corrective action directly;
\ref{sec:rw-hybrid} covers combining measurement and prediction. Section~\ref{sec:rw-eval} then
examines how these methods are evaluated, and argues that the field's reporting conventions make
the individual results impossible to compose into a comparison. That observation is the gap this
thesis addresses.

% -------------------------------------------------------------------------------------
\section{Multirotor dynamics and baseline control}
\label{sec:rw-baseline}

Each vehicle $i$ in a team is modelled as a rigid body,
%
\begin{align}
  \dot{p}_i &= v_i, &
  m\,\dot{v}_i &= -m g e_3 + f_i R_i e_3 + f_{\mathrm{res},i}, \label{eq:rw-trans}\\
  \dot{R}_i &= R_i \hat{\omega}_i, &
  J\dot{\omega}_i &= -\omega_i \times J\omega_i + \tau_i + \tau_{\mathrm{res},i},
  \label{eq:rw-rot}
\end{align}
%
in a world frame with $e_3$ pointing up, where $f_i$ and $\tau_i$ are the collective thrust and
body torque produced by the rotors. The
terms $f_{\mathrm{res},i}$ and $\tau_{\mathrm{res},i}$ collect every force and moment the nominal
model omits. This chapter is concerned with the case in which they are dominated by the
aerodynamic interaction with neighbouring vehicles. The formulation is developed in
Chapter~\ref{ch:model}.

The baseline against which interaction-aware methods are measured is a geometric tracking
controller on $SE(3)$, which tracks a position trajectory and a single heading direction
coordinate-free, avoiding the representation singularities of a local attitude
parameterisation~\cite{lee2010geometrictracking}. Its guarantees are regional rather than global:
exponential stability for initial attitude errors below \SI{90}{\degree}, and almost-global
exponential attractiveness below \SI{180}{\degree}. It is the controller three of the
strategies compared in this thesis build upon, and the one used uncompensated as a reference
condition. The attitude error function is that work's, verified equation-by-equation against its
companion arXiv paper~\cite{lee2010geometricse3}; the torque law implemented here is a
\emph{reduced} form of it, and Chapter~\ref{ch:control} states exactly which terms differ.

A second property of the quadrotor is used throughout: differential flatness. Position and yaw,
together with their derivatives, determine the full state and the inputs, which allows a
feedforward term to be computed directly from a planned trajectory. Tal and
Karaman~\cite{tal2018indi} exploit this to track position and yaw derivatives up to fourth order,
using flatness-derived feedforward for angular rate and angular acceleration in order to follow
jerk and snap. Flatness recurs in Section~\ref{sec:rw-inject} as a structure that a learned
correction can either preserve or destroy.

% -------------------------------------------------------------------------------------
\section{Aerodynamic interaction in multirotor teams}
\label{sec:rw-physics}

The dominant interaction in a vertically stacked formation is the downwash of the upper vehicle
acting on the lower one. Two properties make it difficult to handle. It is strongly nonlinear,
growing sharply as the vehicles close and decaying quickly with lateral offset; and it is
asymmetric, since the vehicle above is affected far less than the one below.

Recent experimental work characterises the effect directly rather than inferring it from tracking
error. Kiran et al.~\cite{kiran2025downwash} measure forces \emph{and} torques with load cells on
mounted Crazyflies whose roll and yaw are locked and whose motors are driven at a constant hover
voltage, and use particle image velocimetry in a centre-plane light sheet to resolve the wake, for
a single quadrotor and for a pair; the dynamic cases traverse one vehicle on a rail rather than
flying it. Gielis et al.~\cite{gielis2023aggregate} measure
six-degree-of-freedom exogenous forces on a multirotor held on a purpose-built load stand at hover
thrust, while up to three further vehicles fly freely around it, and report strong nonlinearities
in how those interactions aggregate. The vehicle experiencing the disturbance is therefore
constrained rather than in free flight. Both measure the disturbance as a transducer load rather
than inferring it from a flight controller's residual --- actuation is still present in each, at
hover thrust --- which makes them, \emph{in our reading}, the natural ground truth against which a
learned model should be checked. Neither paper makes that claim, and neither validates a learned
model of ours or anyone else's.

Simplified physical models exist. Bauersfeld et al.~\cite{bauersfeld2024airflow} analyse roughly
\SI{16}{\hour} of hover data across six platforms spanning \SIrange{0.23}{6.3}{\kilo\gram}, and
show that once the individual propeller wakes have merged the mean flow is well approximated by a
turbulent jet --- beyond about \num{2.5} motor-to-motor distances $l$ below the vehicle, a length
scale the abstract renders as ``drone-diameters''. That gives a computationally light alternative
to computational fluid dynamics. What it estimates is the \emph{time-averaged magnitude} of the
induced flow \emph{in hover}, not an instantaneous six-degree-of-freedom wrench on a follower. For
the near field the authors report no closed-form first-principles model and point to computational
fluid dynamics instead. Formation flight operates in exactly that near field --- our observation,
not theirs --- which is the practical argument for learning the residual from data rather than
deriving it.

% -------------------------------------------------------------------------------------
\section{Reactive compensation: measuring the residual online}
\label{sec:rw-reactive}

The first family estimates the residual from its effect and requires no prior model of the
interaction. Incremental nonlinear dynamic inversion (INDI) compares measured acceleration with the
acceleration the nominal model predicts for the actuation actually applied, and attributes the
difference to unmodelled forces, which it then cancels incrementally.

Smeur et al.~\cite{smeur2017cascadedindi} cascade an inner angular-acceleration INDI loop with an
outer linear-acceleration INDI loop on a single Parrot Bebop running Paparazzi at
\SI{512}{\hertz}; the position and velocity loops feeding the acceleration reference remain
proportional-derivative rather than a third incremental stage. They demonstrate disturbance
rejection indoors by flying in and out of a \SI{10}{\metre\per\second} open-jet exhaust every
\SI{14}{\second}. In two \emph{separate} experiments they show the method does not depend on
frequent position updates: indoors with motion capture deliberately downsampled to \SI{4}{\hertz},
and outdoors taking off and holding position on a consumer GNSS module in roughly
\SI{5.1}{\metre\per\second} wind. Tal and Karaman~\cite{tal2018indi} use INDI for robust tracking of
linear and angular accelerations in aggressive flight on a single custom \SI{609}{\gram} vehicle,
and state the property that defines the family: no parametric model of the aerodynamic effects is
required. That is narrower than it sounds --- mass, inertia, motor time constant, control
effectiveness and the throttle map are all still identified; it is the \emph{drag} model that is
dispensed with.

The same work also exposes the family's cost. Tracking snap requires direct control of body torque,
which those authors achieve using closed-loop motor speed control based on optical encoders on
their own vehicle's motors~\cite{tal2018indi}; encoders are what \emph{that} snap-and-torque path
required, not a general prerequisite for INDI. Reconstructing the applied wrench well enough to
difference it against measurement does, however, generally require some form of rotor-speed
feedback --- Smeur et al.\ likewise use motor-speed
measurement~\cite{smeur2017cascadedindi} --- which is additional hardware or firmware support that
must be fitted, calibrated and maintained on every vehicle. A second limitation is structural rather than
practical: the estimate is delayed by the sensing and filtering chain, and it exists only once the
disturbance has already perturbed the vehicle. A reactive controller cannot, even in principle, act
before the disturbance arrives.

That the sensing requirement is not fundamental has since been shown. Cobo-Briesewitz et
al.~\cite{cobobriesewitz2026lindi} demonstrate that a neural network can generate smooth
approximations of INDI's residual estimate without rotor-speed measurement, which is discussed in Section~\ref{sec:rw-hybrid}.

% -------------------------------------------------------------------------------------
\section{Predictive compensation: learning the residual from relative state}
\label{sec:rw-predictive}

The second family predicts the residual before it acts, from the relative states of the
neighbours, using a model learned offline from flight data. Because the interaction is a function
of the team's relative configuration, a model with the right inputs can act ahead of the
disturbance rather than after it.

Shi et al.~\cite{shi2020neuralswarm} combine a nominal dynamics model with a regularised
permutation-invariant deep neural network that learns multi-vehicle interaction, and use the
learned model as a feed-forward term in a decentralised nonlinear tracking position controller;
geometric \(SE(3)\) control appears only as an optional attitude inner loop, and the learned
force itself is the vertical component \(f_{a,z}\), evaluated onboard the STM32. Training
separate networks on data from two, three and four vehicles respectively, they report worst-case
height error reduced by a factor of two or better, and up to four times smaller in the most
favourable case. In a two-vehicle swapping task the baseline reaches \SI{9.4}{\centi\metre};
the network trained on \emph{four} vehicles brings this to \SI{2.4}{\centi\metre}, while the
network trained on two vehicles reaches \SI{2.7}{\centi\metre}. The swap experiments use vertical
neighbour spacings of \SIrange{0.20}{0.25}{\metre}, with \SI{25}{\centi\metre} stated explicitly
for the three-vehicle case. Generalisation is directional rather than universal: networks trained
on three or four vehicles improved five-vehicle flight, whereas a network trained only on two
vehicles did \emph{not} generalise, performing worse than the baseline on three- and four-vehicle
tests.

Neural-Swarm2~\cite{shi2022neuralswarm2} extends this to heterogeneous teams using heterogeneous
deep sets, whose permutation-invariant architecture accepts any number of neighbours. That is a
structural property rather than an accuracy claim: the networks are trained on at most three
vehicles (at most two neighbours), and the paper describes prediction at sixteen robots as
relatively less accurate. Two design choices are worth noting. Spectral normalisation is applied,
which the abstract describes as yielding stability and generalisation guarantees; in the body this
rests on Lipschitz bounds together with empirical generalisation and a bounded-residual assumption
used in the proofs, rather than a formal guarantee on unseen data. The learned residual is used in
interaction-aware motion planning as well as in tracking control. They report flying a
heterogeneous team of sixteen robots at a minimum vertical distance of \SI{24}{\centi\metre}, and,
separately, up to a threefold reduction in worst-case tracking error against a baseline nonlinear
tracking controller with the interaction force set to zero. The two figures are reported in the
abstract and we do not assume they describe the same experiment: on the sixteen-robot flight
itself the body reports the maximum \(z\)-error of a small vehicle falling from \SI{15}{\centi\metre}
to \SI{10}{\centi\metre}, a factor of about \num{1.5}.

Subsequent work has focused on making these models cheaper to obtain rather than larger. Smith et
al.~\cite{smith2023so2equivariant} exploit the rotational symmetry of the problem with an
$SO(2)$-equivariant model, and report that five minutes of real-world flight data attains a lower
validation error than learning baselines trained on more than fifteen --- a data-efficiency result
on prediction accuracy rather than a closed-loop tracking one. Deployed online within an optimal
feedback (LQR) controller on two identical vehicles, the model reduces three-dimensional tracking
error by \SI{36}{\percent} and vertical tracking error by \SI{56}{\percent} on
average.\footnote{The paper is internally inconsistent on whether the \SI{36}{\percent} figure is
three-dimensional or lateral: we follow the abstract and the translation-result body, which give
three-dimensional and vertical. Its contributions list instead calls the \SI{36}{\percent} figure
lateral.} On a lateral transect flown beneath the hovering vehicle they report mean
three-dimensional position error falling from \SI{0.154}{\metre} to \SI{0.098}{\metre}, close to
\SI{36}{\percent}, with vertical improvements of \SI{55}{\percent} and \SI{49}{\percent}. The
platform is not a Crazyflie but a custom quadrotor of roughly \SI{0.26}{\metre} span and
\SI{0.7}{\kilo\gram}, tracked by OptiTrack and controlled from a Raspberry~Pi at \SI{50}{\hertz}
with a \SI{45}{\hertz} network update. The closest \emph{training} stage is about
\SI{0.5}{\metre} --- some two body-lengths --- while the evaluation manoeuvres are flown at
\SIrange{0.6}{0.8}{\metre}, together with a lemniscate flown beneath the hovering vehicle at
\(z = \SI{-2.5}{\metre}\). The authors
attribute the efficiency to exploiting the problem's latent geometry rather than learning it from
data.

\subsection{The superposition assumption}
\label{sec:rw-superposition}

The deep-sets architectures above aggregate per-neighbour terms by summation, which commits them
to a form of superposition. The literature is explicit that this is an approximation.
Shi et al.~\cite{shi2020neuralswarm} observe that with more than two vehicles the aerodynamic
effect is not a simple superposition of each pair. Gielis et al.~\cite{gielis2023aggregate} make
the same point as their motivation --- an accurate model of the downwash from one vehicle is
unlikely to suffice for predicting the aggregate from several. Their abstract frames this as
modelling the formation as a graph; the implementation learns the aggregation function with a
sum-pooling deep-sets predictor, contrasted against a linear-summation baseline. Generalisation to
a fourth vehicle is claimed at abstract level, while the reported tables stop at three neighbours
and the conclusion is pessimistic about extrapolating further, so we do not treat it as a
demonstrated held-out result.

Their measurements refine the picture rather than simply refuting summation. They report that
interactions fall into three regimes depending on the formation and its location in the state
space: a chaotic, noise-dominated regime; a linear regime \emph{well estimated by linear
summation}; and a nonlinear regime requiring a nonlinear model. Superposition is therefore not
uniformly wrong but regime-dependent, which makes the practical question one of how much accuracy
is lost, and where. This is their own load-stand result, obtained independently of Shi et al.'s
free-flight observation cited above.

Summed architectures nevertheless remain in wide use, because a controller does not require an
exact model, only one whose remaining residual is smaller than the one it removes. What is not
reported is the size of the penalty that the simplification carries for a deployed model, which is
the question this thesis takes up in the multi-robot results chapter.

% -------------------------------------------------------------------------------------
\section{Injecting the prediction: the controller matters}
\label{sec:rw-inject}

A predicted residual is inert until it enters a control law, and the manner of entry differs by
controller family. Treating ``predictive compensation'' as a single method would obscure
differences that are as large as those between the families themselves.

\paragraph{Feedforward into a geometric controller.}
The simplest injection adds the predicted force to the desired force of the baseline controller,
leaving its structure unchanged. This is the form used with the permutation-invariant models of
Section~\ref{sec:rw-predictive}, whose baseline is a nonlinear tracking
controller~\cite{shi2020neuralswarm,shi2022neuralswarm2}.

\paragraph{Feedback linearisation with a flatness-preserving residual.}
A learned correction added to a nominal model can destroy differential flatness, and with it the
applicability of flatness-based planning and control. Yang et al.~\cite{yang2025flatnessresiduals}
show that residuals with a lower-triangular structure preserve both the local flatness of the system
and its original flat output for systems whose nominal model admits a pure-feedback form, and give
a constructive procedure for recovering the diffeomorphism of the augmented system. The residual
there is a \emph{structured correction inside the state derivative} rather than a free
six-degree-of-freedom wrench --- in their worked example it appears only in the velocity equation.
Their validation is in simulation on a planar, two-dimensional quadrotor; there is no hardware, no
three-dimensional vehicle and no multi-robot downwash experiment, and they list a physical quadrotor
as future work. Hsieh et
al.~\cite{hsieh2026flatness} apply the idea to tight formation flight, learning a physics-informed
residual that keeps the joint multi-quadrotor system differentially flat and using the preserved
structure to design a feedback-linearising controller that cancels the disturbance by feedforward.
Their abstract reports a \SI{31}{\percent} reduction in average tracking error against nominal
baselines; the body does not restate that aggregate, reporting instead, on hardware, \SI{24}{\percent}
RMSE and \SI{30}{\percent} maximum-\(z\) improvement for ROM-FBL over nominal FBL, and a further
\SI{29}{\percent} RMSE and \SI{43}{\percent} maximum-\(z\) improvement for learned FBL over
ROM-FBL. On hardware they match the approximately \SI{5}{\centi\metre} RMSE and \SI{10}{\centi\metre}
maximum-\(z\) error \emph{published} by Chee et al.~\cite{chee2024knodedwmpc}, rather than
comparing the two methods on one rig; in their simulation study the flatness-based controller's
maximum-\(z\) error can be worse than that of nonlinear model predictive control. The \num{20}-fold
reduction in solve time (about \SI{1}{\milli\second} against \SI{27}{\milli\second} at equal
residual) is a \emph{simulation} figure; the introduction separately quotes a sevenfold compute
reduction, and the abstract an order of magnitude. Training uses \SI{28}{\second} of hardware data
and \SI{18}{\second} in simulation. Two qualifications matter for our purposes: the hardware
experiments use two Crazyflie~2.1 quadrotors with thrust-upgrade bundles --- very nearly the
platform used in this thesis --- and the controller runs \emph{off-board}, on an Intel~i5 laptop
streaming commands over a Crazyradio link at \SI{400}{\hertz}, rather than on the vehicle. Their
\SI{5}{\milli\second} loop is therefore an abstract-level, off-board figure, and the claim it
supports is that the method is cheap enough to be worth porting, not that it has been shown to run
onboard.

On hardware they also compensate residual \emph{forces} only --- commanding thrust and angular
velocity, and delegating residual torques to the vehicle's own rate controller --- the same scoping
decision taken in Chapter~\ref{ch:model}. Their simulation study does still model residual torques.

\paragraph{Inside a model predictive control horizon.}
The residual can instead form part of the prediction model that an optimiser reasons over, so that
the controller plans against a disturbance it expects to meet rather than cancelling the one
present now. Chee et al.~\cite{chee2024knodedwmpc} embed a mixed-expert KNODE-DW model --- first
principles plus a deliberately lightweight neural ODE component --- inside the prediction
constraint of a nonlinear MPC, rather than adding it as an instantaneous feedforward term, and
evaluate it in both simulation and hardware. The hardware experiments use two Crazyflie~2.1
vehicles in stacked formation, the lower one carrying a thrust-upgrade bundle, flown with average
separations below \num{1.5} body lengths --- one case averaging \(\Delta z = \SI{0.12}{\metre}\),
about \num{1.2} body lengths --- with commanded separations from \SIrange{0.2}{0.4}{\metre} down to
a very tight \SI{0.08}{\metre}. Against a nominal MPC baseline they report average RMSE better by
\SI{40.1}{\percent} and maximum \(z\)-error smaller by \SI{57.5}{\percent}, from \SI{46}{\second}
of training data in total (\SI{18}{\second} simulated and \SI{28}{\second} on hardware). As with
the flatness-based work above, the controller runs \emph{off-board} --- an Intel~i5 laptop
commanding over Crazyradio at about \SI{400}{\hertz} with Vicon at about \SI{120}{\hertz}, the
onboard firmware retaining PID control of thrust and body rates.
Li et al.~\cite{li2023ndpnmpc} take a related route, training a multilayer perceptron --- spectrally
normalised on its weight matrices each epoch --- to predict the downwash force from the neighbour's
relative position and velocity, then injecting the predicted force sequence into the nonlinear MPC's
discrete prediction model as optimiser parameters, so that the lower vehicle can prepare in advance
from the neighbour's planned trajectory rather than from its current state alone. Their evaluation
is on hardware with two similar custom quadrotors, the optimisation running onboard a Jetson~TX2~NX
with pose from off-board motion capture. Hsieh et
al.~\cite{hsieh2025onlineadaptation} extend the approach by inserting an $\mathcal{L}_1$
residual-acceleration estimate into the KNODE-DW prediction model, and evaluate it on \emph{three}
Crazyflie~2.1 vehicles in V-stack and I-stack formations --- including a vertically aligned stack
with pairwise separations of \SI{0.2}{\metre}, about two body-lengths. Only the $\mathcal{L}_1$
estimate adapts online; the neural component reuses the training data of the work above and is not
retrained in flight. The evaluation is hardware-only, again off-board, with the vehicles driven
from separate machines and commanded in thrust and body rates. They report that \emph{``the
$\mathcal{L}_1$ adaptive module compensates for unmodeled disturbances most effectively when paired
with an accurate dynamics model''} --- in our reading, that online adaptation complements a good
residual model rather than substituting for one.

\paragraph{Inside an optimal feedback controller for a docking task.}
A fourth instance is worth separating because its task differs in kind. Shankar et
al.~\cite{shankar2023docking} embed a learned lumped downwash force online in an infinite-horizon
LQR acting on a feedforward-linearised follower, beneath an inner attitude loop, and use it for
vertical docking: a leader--follower manoeuvre in which the vehicles close to within one or two
body-lengths, so that the disturbance is not incidental to the task but constitutive of it. The
platforms are two custom quadrotors of roughly \SI{0.27}{\metre} diagonal and \SI{0.7}{\kilo\gram},
not the Crazyflie class, and the free-flight dock is performed at a vertical offset of about
\SI{0.36}{\metre} --- close to one body-length --- with a curriculum descending to around
\SI{0.5}{\metre}. The model itself is the $SO(2)$-equivariant one of their prior work, trained
offline from about five minutes of data and queried online on an onboard Raspberry~Pi, with pose
from motion capture; the evaluation is hardware only. Reported accuracy depends on the manoeuvre:
the abstract gives terminal error below \SI{0.06}{\metre}, the body reports terminal vertical error
typically below \SI{0.1}{\metre} for a hovering leader, and for a moving leader below
\SI{0.05}{\metre} with the model against more than \SI{0.2}{\metre} without it. Physical docks
succeeded in eight of ten attempts against a hovering leader at \SI{0.36}{\metre}; against a leader
moving at \SI{0.25}{\metre\per\second} success was limited, which the authors attribute to gripper
latency. It is the closest published setting to the tightest
scenarios of Chapter~\ref{ch:setup}, and the regime in which --- on our reading ---
compensation ceases to be an accuracy improvement and becomes a precondition for the manoeuvre.

% -------------------------------------------------------------------------------------
\section{Learning the compensation directly}
\label{sec:rw-rl}

A third stance declines to model the force at all and learns the corrective action instead. Zhang
et al.~\cite{zhang2024proxfly} train a residual reinforcement-learning module on top of a cascaded
controller, producing high-level commands that compensate the disturbance and thrust loss caused by
downwash, and use domain randomisation to relax the requirement for accurate system identification.
They train in simulation and evaluate in both simulation and hardware, on a heterogeneous pair of
regular-sized quadcopters rather than the Crazyflie class. The residual and the high-level
controller run \emph{off-board} at \SI{50}{\hertz}, with the low-level loop onboard at
\SI{500}{\hertz} and motion capture at \SI{200}{\hertz}, so this is not an end-to-end onboard
residual controller.

One property distinguishes this approach sharply from everything in
Section~\ref{sec:rw-predictive}: the policy's observation comprises the tracking error, the
baseline cascaded controller's command and its own previous residual action --- no neighbour state,
and no communication between vehicles~\cite{zhang2024proxfly}. Every relative-state method depends on knowing where the
neighbours are, whether by communication or by onboard sensing. Removing that dependency removes a
bandwidth requirement and a failure mode, at the cost of the ability to anticipate a neighbour the
vehicle cannot perceive.

The approach also gives up the explicit force estimate that the other families produce. That
estimate is what makes a predicted residual comparable against a measured one, and its absence
makes failures correspondingly harder to attribute --- a consideration in experimental design as
much as in deployment.

% -------------------------------------------------------------------------------------
\section{Hybrid reactive and predictive compensation}
\label{sec:rw-hybrid}

Measurement and prediction are complementary in principle. Prediction is available before the
disturbance acts but is valid only over the configurations its training data covered; measurement
is general but delayed. Combining them is therefore natural, and has been done for a single
vehicle.

Cobo-Briesewitz et al.~\cite{cobobriesewitz2026lindi} propose two methods. In the first, LINDI, a
neural network replaces INDI's residual estimate entirely, trained on smoothed data obtained by
spline fitting. In the second, NA-INDI, the total residual is written \(f_a = f_{a,\mathrm{NN}} +
f_{a,\mathrm{INDI}}\): the network predicts a component and the incremental term reasons about the
remaining mismatch. Both are evaluated on a single quadrotor, with and without
a slung payload, on figure-of-eight and circular trajectories --- neither of which was in the
training set.

The headline result concerns sensing rather than the hybrid: a network can generate smooth
approximations of INDI's output without requiring rotor-speed sensors at LINDI's runtime, which
removes the hardware prerequisite identified in Section~\ref{sec:rw-reactive} from that method's
deployment --- though rotor speeds are still needed to produce the training labels, and NA-INDI
retains an incremental branch that uses them. Without a payload, residual compensation of any kind
reduces tracking error by around \SI{40}{\percent} relative to the geometric baseline; with a
payload the reductions are smaller, in the range of \SIrange{19}{24}{\percent}. The hybrid's own
advantage is real but modest: without a payload it gives the best performance, though the authors
describe the margin as small, and with a payload the authors describe it as remaining
\emph{similar} to INDI alone. They also observe a slight drop in performance on the circular
trajectory --- for LINDI, in the no-payload case --- attributing it to that trajectory being out of
distribution relative to the training data, a direct instance of the limitation that defines the
predictive family.

Two features of this setting bear on whether the result transfers. The residual there is
self-induced --- blade flapping, drag, the swing of a payload --- and is predicted from the
vehicle's own state. In a multi-robot team the residual is produced by other vehicles and is a
function of a measurable relative state, which gives the predictive component an input the
single-vehicle setting does not provide. Whether that turns a small margin into a decisive one is
untested, and is one of the questions this thesis addresses.

% -------------------------------------------------------------------------------------
\section{How these methods are evaluated}
\label{sec:rw-eval}

The preceding sections describe a field with strong results. Reported improvements include a
factor of two or better --- up to four --- in worst-case height error~\cite{shi2020neuralswarm},
up to a threefold reduction in worst-case tracking error against a zero-interaction
baseline~\cite{shi2022neuralswarm2}, \SI{36}{\percent} and \SI{56}{\percent} improvements in
three-dimensional and vertical tracking~\cite{smith2023so2equivariant}, \SI{31}{\percent} against
nominal baselines as stated in the abstract~\cite{hsieh2026flatness}, and roughly
\SI{40}{\percent} for residual compensation without a payload~\cite{cobobriesewitz2026lindi}.
Taken individually, each is credible and internally validated.

They cannot, however, be composed. Table~\ref{tab:rw-conditions} lists the conditions under which
they were obtained.

\begin{table}[t]
  \centering
  \caption{Evaluation conditions across the interaction-force literature. Each result is reported
    against that work's own baseline, and no two rows share a platform, a team size, a separation
    and a trajectory class.}
  \label{tab:rw-conditions}
  \small
  \begin{tabular}{@{}llll@{}}
    \toprule
    Work & Team size & Separation & Reported quantity \\
    \midrule
    Shi et al.~\cite{shi2020neuralswarm}        & 2--4 train, 2--5 test & \SIrange{0.20}{0.25}{\metre} vert. & max.\ $|z|$ error \\
    Neural-Swarm2~\cite{shi2022neuralswarm2}    & 16, heterogeneous  & \SI{24}{\centi\metre} vert. & worst-case tracking error\textsuperscript{\dag} \\
    Smith et al.~\cite{smith2023so2equivariant} & 2                  & \SIrange{0.6}{0.8}{\metre} eval.\ (train.\ to $\approx$\SI{0.5}{\metre}) & 3D and vertical tracking \\
    Gielis et al.~\cite{gielis2023aggregate}    & up to 4 total (rig sufferer $+\leq3$ flyers) & test rig, sufferer not in free flight & force-prediction accuracy \\
    Hsieh et al.~\cite{hsieh2026flatness}       & 2                  & \SI{0.3}{\metre} and \SI{0.4}{\metre} vert.\textsuperscript{\ddag} & RMSE and max.\ $z$ error \\
    Hsieh et al.~\cite{hsieh2025onlineadaptation} & 3                & V-stack and I-stack; pairwise \SI{0.2}{\metre} & tracking vs.\ MPC family \\
    Cobo-Briesewitz et al.~\cite{cobobriesewitz2026lindi} & 1 (+ payload) & n/a                 & position tracking error \\
    Zhang et al.~\cite{zhang2024proxfly}        & 2, heterogeneous   & \SI{0.50}{\metre} vert.\ hover/circle; docking approach to \SI{0.10}{\metre}, then release at \SI{0.05}{\metre} & position and attitude error \\
    \bottomrule
  \end{tabular}

  \vspace{0.4em}
  \begin{minipage}{\linewidth}\footnotesize
    \textsuperscript{\dag}~The sixteen-robot flight and the threefold reduction are separate
    reports; on the sixteen-robot flight the body gives a factor of about \num{1.5}.
    \textsuperscript{\ddag}~Hardware tracking comparisons are at commanded
    \(\Delta z = \SI{0.3}{\metre}\) and \SI{0.4}{\metre}; the \SIrange{0.2}{0.5}{\metre} range
    belongs to the simulation and prediction evaluation.
  \end{minipage}
\end{table}

A \SI{36}{\percent} improvement measured on two vehicles in close proximity and a
\SI{31}{\percent} improvement measured in a tight formation are not the same quantity: they differ
in what was flown, on what airframe, at what separation, against what baseline, and under what
metric. Neither can be subtracted from the other, and their ordering carries no information about
which method would prevail if both were flown under the same conditions. The difficulty compounds
where the baselines themselves differ, since an improvement relative to a weaker baseline is
easier to obtain.

This is not a criticism of any individual work. Each is internally valid, and there is no
convention in the field that any of them violated. It is a limitation of the literature taken as a
whole: the results are not designed to compose, and no amount of careful reading assembles them
into a ranking. Establishing one requires an experiment built for the purpose, which is what
distinguishes a comparison from a survey.

\subsection{Comparisons that do exist}
\label{sec:rw-existing-comparisons}

The field does conduct systematic comparisons, and two recent ones bound what remains open.

Gu et al.~\cite{gu2025comparative} benchmark four online-learning uncertainty estimators ---
extreme learning machines, radial-basis-function networks, echo state networks and Gaussian
processes --- inside a single modular framework in which a nominal controller is paired with a
swappable residual-learning module. They treat the uncertainty as a lumped additive wrench
\((f_a, \tau_a)\) recovered by inverting the discrete Newton--Euler dynamics --- analogous to the
residual acceleration used here, though not identical to it a priori. They evaluate estimation accuracy, tracking
performance and computational cost, and reach the kind of conclusion a comparison should: not a
winner, but a set of trade-offs, with Gaussian processes most accurate yet limited by cubic
training complexity, the lightweight parametric models fast but sensitive to initialisation, and
echo state networks the most balanced.

Two things separate it from the present work. Its subject is a \emph{single} quadrotor under
generic uncertainty, and all four methods learn \emph{online}; it therefore compares estimators
\emph{within} one family rather than the families of Sections~\ref{sec:rw-reactive}
to~\ref{sec:rw-hybrid}. And its evaluation is in simulation throughout: the reported hardware
component is the learning modules executing on an embedded computer in the loop with a simulator,
which measures latency rather than flight, and the authors list validation on a real quadrotor as
future work.

Its approach to fairness is nonetheless directly relevant, and is the opposite of the one taken
here. They tune the parametric models --- the extreme learning machine, the radial-basis-function
network and the echo state network --- to their own optima by Bayesian optimisation under
otherwise consistent conditions, so that each operates near its best achievable performance; this
thesis freezes a shared gain set instead. Both answer the same objection --- that an observed difference
might reflect tuning effort rather than method --- and Section~\ref{sec:ctrl-fairness} returns to
why the choice is genuinely open.

Abro et al.~\cite{abro2025fractional} compare fractional-order PID, fractional-order LQR and
integral state feedback on close-proximity flight, explicitly motivated by downwash and by the
excessive safety distances that neglecting it imposes.\footnote{Cited for the methods compared and
for that motivation; the abstract's quantitative improvement figure and its claim about larger
swarms are not used here.} This is a controller comparison in our setting: multiple vehicles, and
downwash-aware. Of the three, only integral state feedback is classical; the fractional-order laws
are not. All three are nonetheless feedback regulators: they do not present a learned or
first-principles predictor of the interaction force in the sense of the works above, so the
reactive, predictive and hybrid families remain uncompared against one another.

Together these delimit the gap precisely. It is not that close-proximity control is uncompared, nor
that the field lacks the methodology. It is that the specific families defined by \emph{how they
obtain information about the interaction force} --- measure it, predict it, or both --- have not
been placed under common conditions.

Such an experiment has requirements that follow directly from the failure modes above: a single
airframe; identical trajectories and formations across all methods; gains frozen before data
collection and shared wherever the architecture permits, so that a difference cannot be a tuning
artefact; a single definition of the residual, measured identically under every controller
including those that do not use it; and a null condition in which no interaction is present, to
establish the measurement floor and so distinguish a real effect from noise. Chapter~\ref{ch:setup}
describes the resulting design.

% -------------------------------------------------------------------------------------
\section{Summary and gap}
\label{sec:rw-gap}

A controller can obtain information about the interaction force in two ways. It can measure the
residual online, which requires no prior model and generalises to disturbances never seen, but is
delayed and --- in its established form --- depends on rotor-speed sensing. Or it can predict the
residual from the relative states of its neighbours using a model learned offline, which acts ahead
of the disturbance but is valid only where the training data reached. A third approach learns the
corrective action directly, forgoing an explicit force estimate and, in at least one instance,
inter-vehicle communication as well. Representative methods exist for each, and each is
individually validated.

Three gaps follow from this survey.

First, and principally, to the best of our knowledge the reactive, predictive and hybrid families
have not been evaluated against one another under identical conditions. Systematic comparisons in
adjacent territory exist --- of online-learning estimators on a single vehicle, and of classical
regulators for close-proximity swarms (Section~\ref{sec:rw-existing-comparisons}) --- but none
spans the families distinguished by how they obtain the interaction force. The individual results are strong; what the field lacks is not strong results
but comparable ones, and its reporting conventions --- improvements against each work's own
baseline, on its own platform, at its own separations --- make the comparison impossible to
assemble after the fact.

Second, the hybrid of measurement and prediction has been demonstrated only for a single vehicle,
where the residual is self-induced and its advantage was described as small without a payload and
as leaving performance similar to the reactive method with one~\cite{cobobriesewitz2026lindi}. The
multi-robot case, in which the residual is larger and is
a function of a measurable relative state, is where the predictive component has the most to
contribute and remains untested.

Third, the summation used to aggregate per-neighbour predictions is known to be an
approximation~\cite{shi2020neuralswarm,gielis2023aggregate}, yet the magnitude of the penalty it
carries for a deployed model, evaluated on teams larger than those it was trained on, is not
reported.

\todo[color=green!25]{\textbf{Provenance, not pending work.} Supported by two audits: an arXiv keyword sweep, and a
cross-database sweep over OpenAlex, Crossref and DBLP including a citation-graph scan of 157 works
citing the Neural-Swarm line. Both found adjacent comparisons but none spanning these families.
Neither is exhaustive; keep the ``to the best of our knowledge'' phrasing.}

Chapter~\ref{ch:model} formalises the residual wrench these methods estimate or predict, and
Chapter~\ref{ch:control} describes the seven strategies compared in this thesis and the point at
which each injects its knowledge of the interaction into the control law.
